\section{Question, scope, and relation to the existing programme}
\label{sec:context}

The immediate problem is to make complete normal-surface search less
expensive without weakening its mathematical meaning.  A fast test of
one supplied surface, a fast search for one favourable linear objective,
and a fast enumeration of every relevant vertex surface are different
operations.  An improvement in the first does not automatically improve
the worst-case complexity of the third, and none by itself bounds the
number of triangulations or sectors that a recognizer must visit.

The present continuation targets a specific expense exposed by the
preceding report: once matching has reduced a sector to a small
quadrilateral subspace, the existing general arrangement routine still
intersects equality hyperplanes that never change the canonical lift.
In dimension three, the normalized feasible set is a polygon.  The
relevant geometry is therefore the overlay of the actual planar regions
on which each triangle potential attains its minimum.  Enumerating those
regions gives both a stronger output bound and a direct certificate that
the whole supplied sector has been covered.

The second expense is earlier in the pipeline.  The original constructor
materializes a length-\(k\) quadrilateral label for every face-corner
equation, although all but \(O(k)\) labels vanish.  Compiling the
source incidences once and postponing dense labels until after this
test removes that \(tk\) allocation pattern.  These two changes address
separate stages, and the experiments time them separately.

\subsection{The object actually being enumerated}

Let \(T\) be a finite compact triangulation with \(t\) tetrahedra.
Normal coordinates use four triangle counts and three quadrilateral
counts in each tetrahedron.  A \emph{supplied sector} chooses at most
one allowed quadrilateral type per tetrahedron and sets all other
quadrilateral coordinates to zero.  Its number of allowed types is
\(k\le t\).  Allowed coordinates may vanish: a sector contains all its
coordinate faces.  The remaining matching equations and nonnegativity
inequalities define a rational polyhedral cone in standard coordinates.
We enumerate its extreme rays with nonzero quadrilateral part, in
primitive integer normalization.  The excluded rays are pure vertex
links.

This is a complete list of the nonlink \emph{standard} rays of the
sector.  It is not merely a list of extreme rays in quadrilateral
coordinates.  Canonical lifting is positively homogeneous but
piecewise linear, so a quadrilateral direction in the interior of a
quadrilateral face can lift to a standard extreme ray.  This distinction
is established in the classical conversion theory~\cite{Burton2009}
and is essential here: the geometric family in
Section~\ref{fib:family} has four essential standard disc rays but only
one of them is quadrilateral-extreme.

The abstract cone arguments apply whenever the stated potential model
is available.  The executable modules use the maintained native
validator, whose contract is narrower: a connected orientable finite
3-manifold triangulation with one torus boundary component.  They do
not implement ideal-vertex or arbitrary boundary-pattern input.
The matching dimension \(d\) is the dimension of the rational
quadrilateral matching kernel before nonnegativity.  The implemented
eligibility test is \(d\le3\).  Nonnegativity may reduce the feasible
dimension further; the empty, point, segment, and polygon cases all
occur in genuine audited sources.

\subsection{What is new in this continuation}

The main mathematical statements are as follows.

\begin{enumerate}
\item \textbf{Reusable sparse construction.}
Theorem~\ref{sp:constructor} proves a pre-elimination work and storage
bound while preserving the incumbent kernel's exact combinatorial
representatives, potential gauge, cycle rows, and rational basis.
\item \textbf{Complete planar enumeration.}
Theorem~\ref{pl:quadratic-bound} gives a bound in terms of the
\emph{active} minimum cells, rather than all pairwise potential
equalities.  Proposition~\ref{pl:operations} gives an exact polynomial
construction and accounts for writing the full normal vectors.
\item \textbf{Independent completeness replay.}
Theorems~\ref{cv:area-theorem} and~\ref{cv:complete-replay} show why
dominance and exact total area suffice to certify maximal minimum
cells, not just a partial collection of valid rays.
\item \textbf{An infinite geometric test family.}
Section~\ref{fib:family} solves the entire cone and topology of the
two-cap Fibonacci family, including the exact growth of the old
arrangement's rejected directions.  Section~\ref{cf:grid} separately
shows quadratic sharpness for abstract potential cones.
\end{enumerate}

The canonical-lift characterization itself is not a new claim.
It is proved here to fix every degeneracy needed by the algorithm,
with attribution to both the earlier report and the classical
conversion framework.  Planar minimum diagrams are clipped power
diagrams; their computational geometry and the quadratic complexity of
three-dimensional normal-fan overlays have substantial prior
literature~\cite{Aurenhammer1987,Weibel2010,Fogel2009}.
Our contribution is the explicit complete normal-sector algorithm,
the parameter-sensitive boundary count, its independent certificate,
the exact implementation, and the geometric family.  We make no claim
to a first general theorem about planar subdivisions.

\subsection{Repository lineage}

All code and incoming reports were read against the pinned commit
\begin{center}\small\ttfamily\basecommit\end{center}
The immediately preceding minimum-envelope
package proves complete enumeration for \(d\le2\) and the general
common-cell characterization~\cite{MinimumEnvelopes}.  Its optimized
rank-two implementation is retained unchanged as a performance
control.  The other recent incoming packages cover different parts
of the programme: certified active-face reduction and batched local
descent; bounded trace search; boundary assembly kernels; rooted disc
components; and affine orbit profiles
~\cite{ActiveFaces,TraceSearch,CycleKernels,RootedDiscs,AffineProfiles}.
Their bounds are not imported as global coverage or termination
theorems for this algorithm.
Likewise, global admissible-face bounds proved for closed
triangulations~\cite{Burton2011} are not transferred here to
finite exteriors with boundary.

In particular, the phrase \emph{cycle kernel} is potentially misleading
across the reports.  Here it means linear consistency equations around
cycles in the graph of triangle-coordinate differences.  The
cycle-envelope incoming report studies a boundary assembly language.
These are distinct objects with distinct parameters.

The current maintained source already includes a native, certified
diagram-to-exterior route; its pulling subdivision constructs
\(20\max(1,n)\) tetrahedra from an \(n\)-crossing diagram.  We do not
describe source provenance as an absent component.  The remaining
question is how to use such a source, complete local searches, and
certified topology transformations with a sufficiently small
\emph{total} search bound.  The present modules accept a supplied
triangulation and do not independently authenticate an accompanying
knot diagram.

\subsection{The asymptotic target}

We use \emph{quasi-polynomial} for time
\(2^{(\log n)^{O(1)}}\).  The more specific targets
\(n^{O(\log n)}\) and \(n^{O((\log n)^2)}\) are different.
Lackenby's March 2021 handout announces
\[
2^{O((\log n)^3)}=n^{O((\log n)^2)}
\]
for unknot recognition~\cite{Lackenby2021}.  It describes a hierarchy
strategy combining controlled surface complexity, compressed
representations, multisurfaces, and shortened hierarchy depth.
The July 2026 manuscript gives detailed algorithms and certification
theorems for incompressible surfaces and essential hierarchies, while
explicitly leaving speed-up for further work~\cite{Lackenby2026}.
We therefore use those sources as structural guidance, without
attributing a fully implemented quasi-polynomial analysis to that
manuscript or to the present code.

For comparison, Burton and Ozlen's branching algorithm supplies a
rigorous exponential worst-case analysis and strong experimental
behaviour~\cite{BurtonOzlen2012}.  Improving one subroutine here neither
changes that theorem nor establishes a polynomial empirical law for
complete recognition.  The statements below are local theorems and
reproducible stage measurements; Section~\ref{sec:research} specifies
what would be needed to compose them into a global bound.
